\documentclass[11pt,a4paper]{article}
\usepackage[margin=24mm]{geometry}
\usepackage{amsmath,amssymb,booktabs,microtype}
\usepackage[hidelinks]{hyperref}
\newcommand{\dd}{\mathrm d}
\newcommand{\id}{\mathbb I}
\newcommand{\tr}{\operatorname{tr}}
\newcommand{\ket}[1]{|#1\rangle}
\newcommand{\bra}[1]{\langle#1|}
\title{Route G-R1: A Finite Relational Quantum Clock}
\author{Conditional evolution and the limits of energy-density-only time}
\date{23 September 2026}
\begin{document}
\maketitle
\begin{abstract}
The user-selected mainline now concerns relational time, rather than
motion in the spacetime assumed by the earlier KK model. We give a
finite Page--Wootters-type construction: both local Hamiltonians are
nonnegative, a stationary joint state encodes exact conditional
Schr\"odinger evolution, and the clock is a specified covariant
measurement. An exact coherent/dephased pair has the same complete
energy statistics but different conditional histories. Thus an ordinary
energy density alone cannot specify this relational evolution.
Neither the construction nor the cited optical experiment proves the
source document's density--time law. Quantum mechanics, the subsystem
split, spectra, constraint and joint state remain declared inputs.
\end{abstract}

\section{What changes, and what is assumed}
Chen's sections 1.2, 1.8 and 1.9 associate time with energy information,
its density and event quantities \cite{chen}. We no longer identify that
proposal with the pre-existing Minkowski time of G1 or the flight time
of G2-T. Those calculations remain comparison models, with their
conditional results intact. In particular, no KK radius, negative mode
label or detector source is imported into this construction.

Take $\mathcal H_C\otimes\mathcal H_S$, each of dimension $d\geq2$,
with fixed energy bases $\ket n_C,\ket n_S$, $0\leq n<d$.
Choose one positive energy gap $\epsilon$, not a fitted particle mass:
\begin{align}
 H_C&=\epsilon\sum_{n=0}^{d-1}(d-1-n)\ket n\bra n,&
 H_S&=\epsilon\sum_{n=0}^{d-1}n\ket n\bra n,\\
 E_*&=(d-1)\epsilon,&
 K&=H_C\otimes\id+\id\otimes H_S-E_*\id .
\end{align}
Both local energies are nonnegative. For any $\sum_n|a_n|^2=1$,
\begin{equation}\label{eq:global}
 \ket\Psi=\sum_n a_n\ket n_C\ket n_S,\qquad K\ket\Psi=0 .
\end{equation}
The constraint, not a negative physical energy, subtracts the common
total energy. The dimensionless transformation
$U(\alpha)=\exp(-i\alpha K/\epsilon)$ leaves this state invariant.
$\alpha$ labels a symmetry orbit, not an assumed external physical
elapsed time. This is a conditional quantum construction in the
Page--Wootters tradition \cite{moreva,quantumtime}, not a derivation
of the Born rule, the tensor split or the Hamiltonians from nothing.

\newpage
\section{Clock events and exact conditional evolution}
Specify the clock dial before predicting any system outcome:
\begin{equation}\label{eq:clock}
 \ket\varphi_C=\frac1{\sqrt d}\sum_{n=0}^{d-1}
                  e^{in\varphi}\ket n_C,\qquad
 \varphi_k=\frac{2\pi k}{d},\quad 0\leq k<d.
\end{equation}
The finite geometric sum gives
$\langle\varphi_k|\varphi_l\rangle=\delta_{kl}$ and
$\sum_k\ket{\varphi_k}\bra{\varphi_k}=\id_C$.
For continuously labeled readings, instead use the POVM
\begin{equation}
 F_C(\dd\varphi)=\frac d{2\pi}\ket\varphi\bra\varphi\,\dd\varphi,
 \qquad \int_0^{2\pi}F_C(\dd\varphi)=\id_C .
\end{equation}
Continuous dial states are \emph{not} mutually orthogonal. A POVM with
continuous labels is not infinite clock precision. Its covariance is
\begin{equation}
 e^{-iH_C\alpha/\epsilon}\ket\varphi
       =e^{-i(d-1)\alpha}\ket{\varphi+\alpha}.
\end{equation}
Here phase is a clock-readout label, not a spatial rotation or an extra
geometrical circle. The grid outcome probability is $p(k)=1/d$;
the continuous outcome density is $p(\varphi)=1/(2\pi)$.

For a system effect $0\leq B\leq\id_S$, the joint Born rule gives
\begin{align}
 \ket{\psi(\varphi)}
    &=\sqrt d\,{}_C\bra\varphi\Psi\rangle
      =\sum_n a_n e^{-in\varphi}\ket n_S,\\
 P(B|\varphi)&=\bra{\psi(\varphi)}B\ket{\psi(\varphi)},\\
 i\epsilon\,\partial_\varphi\ket{\psi(\varphi)}
    &=H_S\ket{\psi(\varphi)}.\label{eq:conditional}
\end{align}
This follows by differentiating the finite sum; no external evolution
of $\ket\Psi$ was used. Once the clock is calibrated in time units,
\begin{equation}
 \tau=\frac{\hbar}{\epsilon}\varphi
 \quad\Longrightarrow\quad
 i\hbar\,\partial_\tau\ket{\psi(\tau)}=H_S\ket{\psi(\tau)} .
\end{equation}
This $\tau$ is a calibrated dial reading within one period, not yet
spacetime proper time. Calibration has not derived seconds from $E_d$.
Equivalently, the exact finite history identity is
\begin{equation}
 \ket\Psi=\frac1{\sqrt d}\sum_{k=0}^{d-1}
         \ket{\varphi_k}_C\ket{\psi(\varphi_k)}_S .
\end{equation}
On the grid $\psi_{k+1}=e^{-2\pi iH_S/(d\epsilon)}\psi_k$ and
$\psi_d=\psi_0$. Thus $k$ is cyclic: an ever-increasing event count
requires additional records. Conditioning different ensemble members
on different dial outcomes is not a proof of sequential measurement
statistics on one repeatedly measured clock.

\newpage
\section{Invariant probabilities and finite-clock limitations}
The fixed total-energy subspace is $\ker K$. We can express the same
joint probabilities with effects that preserve the constraint. For
$E_{kB}=\ket{\varphi_k}\bra{\varphi_k}\otimes B$, define
\begin{equation}
 \overline E_{kB}=\frac1{2\pi}\int_0^{2\pi}
        U(\alpha)E_{kB}U(\alpha)^\dagger\,\dd\alpha .
\end{equation}
Since $K/\epsilon$ has integer eigenvalues, the average removes matrix
elements between unequal eigenvalues. Consequently
\begin{equation}
 [K,\overline E_{kB}]=0,\qquad
 \overline E_{kB}\succeq0,\qquad
 \tr(\rho\,\overline E_{kB})=\tr(\rho E_{kB})
       \quad(\rho\text{ supported in }\ker K).
\end{equation}
Twirling the denominator effect in the same way preserves the
conditional probability too. This finite algebraic averaging is a
projection onto invariant effects, not a distribution of an unobserved
physical time. It is consistent with the relational-observable viewpoint
\cite{trinity}. We certify probabilities, not a measurement instrument
or a history of collapse; different instruments can share one effect.

One must not secretly impose an ideal canonical time operator on this
finite clock. For every finite matrix $T$,
\begin{equation}
 \tr[T,H_C]=0\ne i\hbar d .
\end{equation}
Thus $[T,H_C]=i\hbar\id$ is impossible here. A cyclic covariant POVM
avoids that assertion. Neither an irreversible arrow nor unlimited
duration follows. The local generators and covariance also supply the
chosen ordering structure; we have not deduced all notions of order
from correlations alone.

\subsection*{An interaction is a new constraint problem}
Do not insert an arbitrary density-dependent clock speed into
\eqref{eq:conditional}. If $K$ is replaced by $K+V$ and a new physical
state satisfies $(K+V)\ket\Psi=0$, set
$\widetilde\psi(\varphi)=\sqrt d\,{}_C\bra\varphi\Psi\rangle$.
Clock completeness gives the exact projected equation
\begin{align}
 i\epsilon\partial_\varphi\widetilde\psi(\varphi)
 &=H_S\widetilde\psi(\varphi)
 +\frac d{2\pi}\int_0^{2\pi}
 V(\varphi,\varphi')\widetilde\psi(\varphi')\,\dd\varphi',\\
 V(\varphi,\varphi')&={}_C\bra\varphi V\ket{\varphi'}_C .
\end{align}
The projected vector need not have constant norm, and its normalized
conditional equation need not be the same linear equation. A general
$V$ gives a nonlocal kernel in clock labels, not automatically a local
slowdown law. The old $\ket\Psi$ generally fails the new constraint.
Our code checks the projection identity for a test matrix $V$ but
explicitly does \emph{not} claim an interacting stationary completion.

\newpage
\section{Exact counterexample: energy statistics do not fix a history}
For $d=3$, $a_n=1/\sqrt3$, compare
\begin{equation}
 \rho_{\rm coh}=\ket\Psi\bra\Psi,\qquad
 \rho_{\rm diag}=\frac13\sum_{n=0}^2\ket{n,n}\bra{n,n}.
\end{equation}
Both are supported in $\ker K$. They have identical local reductions
$\rho_C=\rho_S=\id/3$ and identical joint energy probabilities
$P(n_C,n_S)=\delta_{n_Cn_S}/3$. In particular,
\begin{equation}
 \langle H_C\rangle=\langle H_S\rangle=\epsilon,\qquad
 \langle H_C+H_S\rangle=2\epsilon,\qquad
 \operatorname{Var}(H_C+H_S)=0
\end{equation}
for both. All functions of the jointly measured local energies have
the same statistics, not merely the same mean. Even if an external
reference volume $V_0$ were supplied, an ordinary mean energy density
$2\epsilon/V_0$ would be identical.

Yet the specified clock measurement yields
\begin{align}
 \rho_{\rm coh}(S|\varphi)&=\ket{\psi(\varphi)}\bra{\psi(\varphi)},\\
 \rho_{\rm diag}(S|\varphi)&=\id/3 .
\end{align}
For the fixed system projector onto
$\ket+=(\ket0+\ket1+\ket2)/\sqrt3$,
\begin{equation}\label{eq:counter}
 P_{\rm coh}(+|\varphi)=
   \frac{|1+e^{-i\varphi}+e^{-2i\varphi}|^2}{9}
   =\frac{(1+2\cos\varphi)^2}{9},\qquad
 P_{\rm diag}(+|\varphi)=\frac13 .
\end{equation}
\begin{center}
\begin{tabular}{rrr}
\toprule
Clock label $\varphi$ & Coherent probability & Dephased probability\\
\midrule
$0$       & $1$   & $1/3$\\
$\pi/3$   & $4/9$ & $1/3$\\
$2\pi/3$  & $0$   & $1/3$\\
$\pi$     & $1/9$ & $1/3$\\
\bottomrule
\end{tabular}
\end{center}
The missing datum is joint coherence. The diagonal state is still a
legitimate stationary quantum state; we do not claim it has \emph{no time}.
It lacks the changing conditional probabilities of this coherent
example. Partial dephasing gives
$P_\eta=\eta P_{\rm coh}+(1-\eta)/3$, $0\leq\eta\leq1$,
without changing the energy statistics.

This refutes a \emph{density-only identification} of the complete
relational history. It does not refute every possible meaning of the
source's unspecified $E_i,E_d$: those could include more than energy
statistics, but then their additional content and dynamics must be
defined. In particular it does not prove $E_d$ causes time dilation.

\newpage
\section{Interpretation, calibration and the next bounded step}
The closest proposed analogue of $E_i$ is not an energy eigenvalue alone
but the joint state \emph{together with} a specified algebra of events.
This is a candidate dictionary, not an identification derived from
Chen's document. The original event quantity $T$ should be denoted
$N_{\rm evt}$ to avoid confusing it with $\tau$ or flight time.
A dial label modulo $d$ is not yet such a count. $E_d$ has no assigned
operator or measure in G-R1; Hilbert dimension $d$ is neither spatial
dimension nor the source's $R$.

There are two simple guards against artificial slowdown:
\begin{itemize}
\item Replacing $(H_C,H_S,E_*,\epsilon)$ by
$a(H_C,H_S,E_*,\epsilon)$, $a>0$, leaves the physical subspace and
$P(B|\varphi)$ unchanged. It rescales $\tau=\hbar\varphi/\epsilon$.
Without an independent reference this is not an observed relative
clock dilation.
\item Shifting $H_C\mapsto H_C+b\id$ and $E_*\mapsto E_*+b$
leaves $K$ unchanged. Clock projectors and correlations are unchanged
although the clock's mean energy shifts. This is a freedom of this
\emph{nongravitational} model, not a claim about energy offsets in a
specified gravity theory.
\end{itemize}
Therefore a physical density--time proposal must choose a density
operator/measure and an interaction, solve the common constraint, and
predict a comparison between calibrated physical clocks. It must
distinguish changes of labels from changes of correlations.

G-R1 is complete within its finite noninteracting scope. The next
bounded task, G-R2, is one explicit initial measurement record and one
final readout. Check sequential conditional probabilities and what
information survives measurement backreaction; do not treat
single-event conditioning as a ready-made accumulating clock.
G-R3 can then test one specified $E_i/E_d$ candidate with a clock-matter
or clock-clock constraint. No freely fitted slowdown function is
authorized. Space, a Lorentzian metric, universal proper time, mass
generation and the source's more expansive claims remain open.

The optical experiment of Moreva et al.\ demonstrates a conditional-clock
mechanism using prepared photons and laboratory apparatus, not that
laboratory spacetime has been derived from information \cite{moreva}.
Work on sequential measurements and invariant relational observables
guides the next step \cite{quantumtime,trinity}; it is not silently
claimed as already implemented here.

\texttt{code/verify\_gr1\_relational\_clock.py} passes 109 algebra checks,
including a nonuniform four-level case, direct projection, numerical
differentiation, invariant effects, the exact counterexample, and
unchanged comparison-file hashes. These are implementation checks,
not new experimental data or a novelty claim for Page--Wootters.

\begin{thebibliography}{9}\small\raggedright
\bibitem{chen} B. Chen, \emph{Energy spectrum for elementary particles},
\href{https://indico.cern.ch/event/1109513/contributions/4969967/attachments/2480973/4410560/Energy\%20spectrum\%20for\%20elementary\%20particles_final.pdf}
{final five-page document}, dated 23 April 2022, Secs.\ 1.2, 1.8, 1.9.
\bibitem{moreva} E. Moreva et al., \emph{Time from quantum entanglement:
an experimental illustration}, Phys.\ Rev.\ A 89, 052122 (2014),
\href{https://arxiv.org/abs/1310.4691}{arXiv:1310.4691}.
\bibitem{quantumtime} V. Giovannetti, S. Lloyd and L. Maccone,
\emph{Quantum Time}, Phys.\ Rev.\ D 92, 045033 (2015),
\href{https://arxiv.org/abs/1504.04215v3}{arXiv:1504.04215v3}.
\bibitem{trinity} P. A. H\"ohn, A. R. H. Smith and M. P. E. Lock,
\emph{The Trinity of Relational Quantum Dynamics},
Phys.\ Rev.\ D 104, 066001 (2021),
\href{https://arxiv.org/abs/1912.00033v3}{arXiv:1912.00033v3}.
\end{thebibliography}
\end{document}
